\documentclass[CJK]{beamer}
\usepackage{CJKutf8}
\usepackage{beamerthemesplit}
\usetheme{Malmoe}
\useoutertheme[footline=authortitle]{miniframes}
\usepackage{amsmath}
\usepackage{amssymb}
\usepackage{graphicx}
\usepackage{color}
\usepackage{slashed}
\usepackage{simplewick}
\graphicspath{{../figures/}}
\def\be{\begin{equation}}
\def\ee{\nonumber\end{equation}}
\def\bea{\begin{eqnarray}}
\def\eea{\nonumber\end{eqnarray}}
\def\ii{{\dot{\imath}}}
\def\bch{\begin{CJK}{UTF8}{gbsn}}
\def\ech{\end{CJK}}
\def\bex{\begin{minipage}{0.3\textwidth}\includegraphics[width=1in]{jugelizi.png}\end{minipage}\begin{minipage}{0.6\textwidth}}
\def\eex{\end{minipage}}
\def\chtitle#1{\frametitle{\bch#1\ech}}
\def\skipline{{\vskip0.1in}}
\def\skiplines{{\vskip0.2in}}
\def\lagr{{\mathcal{L}}}
\def\hamil{{\mathcal{H}}}
\def\vecv{{\mathbf{v}}}
\def\vecx{{\mathbf{x}}}
\def\veck{{\mathbf{k}}}
\def\vecp{{\mathbf{p}}}
\def\vecn{{\mathbf{n}}}
\def\vecA{{\mathbf{A}}}
\def\vecP{{\mathbf{P}}}
\def\vecsigma{{\mathbf{\sigma}}}
\def\hatJn{{\hat{J_\vecn}}}
\def\hatJx{{\hat{J_x}}}
\def\hatJy{{\hat{J_y}}}
\def\hatJz{{\hat{J_z}}}
\def\hatj#1{\hat{J_{#1}}}
\def\hatphi{{\hat{\phi}}}
\def\hatq{{\hat{q}}}
\def\hatpi{{\hat{\pi}}}
\def\vel{\upsilon}
\def\Dint{{\mathcal{D}}}
\def\adag{{\hat{a}^\dagger}}
\def\bdag{{\hat{b}^\dagger}}
\def\cdag{{\hat{c}^\dagger}}
\def\ddag{{\hat{d}^\dagger}}
\def\hata{{\hat{a}}}
\def\hatb{{\hat{b}}}
\def\hatc{{\hat{c}}}
\def\hatd{{\hat{d}}}
\def\hatN{{\hat{N}}}
\def\hatH{{\hat{H}}}
\def\hatp{{\hat{p}}}
\def\Fup{{F^{\mu\nu}}}
\def\Fdown{{F_{\mu\nu}}}
\def\newl{\nonumber \\}
\def\SIkm{\mathrm{km}}
\def\SIyr{\mathrm{yr}}
\def\SIGyr{\mathrm{Gyr}}
\def\SIeV{\mathrm{eV}}
\def\SIGeV{\mathrm{GeV}}
\def\SIm{\mathrm{m}}
\def\SIcm{\mathrm{cm}}
\def\SIJ{\mathrm{J}}
\def\SIs{\mathrm{s}}
\def\SIkg{\mathrm{kg}}
\def\SIg{\mathrm{g}}
\def\vece{\mathrm{e}}
\def\bmat#1{\left(\begin{array}{#1}}
\def\emat{\end{array}\right)}
\def\bcase#1{\left\{\begin{array}{#1}}
\def\ecase{\end{array}\right.}
\def\calM{{\mathcal{M}}}
\def\calT{{\mathcal{T}}}
\def\calR{{\mathcal{R}}}
\def\barpsi{\bar{\psi}}
\def\baru{\bar{u}}
\def\barv{\bar{\upsilon}}
\def\bmini#1{\begin{minipage}{#1\textwidth}}
\def\emini{\end{minipage}}
\def\qeq{\stackrel{?}{=}}
\def\torder#1{\mathcal{T}\left(#1\right)}
\def\rorder#1{\mathcal{R}\left(#1\right)}


\title{Quantum Field Theory I \\ Homework 4 solution}
  \author{}
  \date{}


\begin{document}

\begin{frame}
 
\begin{center}
\begin{Large}
\bch
量子场论 I 

{\vskip 0.3in}

第四次课后作业参考答案
\skipline
\skipline

如发现参考答案有错误请不吝告知（微信zhiqihuang或邮箱huangzhq25@sysu.edu.cn)
\ech
\end{Large}
\end{center}

\vskip 0.2in

\bch
课件下载
\ech
https://github.com/zqhuang/SYSU\_QFTI

\end{frame}

\begin{frame}
\chtitle{第1题：题目和思路}
\bch
题目：对两个矢量$A^\mu$, $B^\mu$，证明
$$\mathrm{Tr}(\slashed{A}\slashed{B}) = 4A^\mu B_\mu$$
其中$\mathrm{Tr}$表示矩阵求迹。


\skipline
思路：利用矩阵乘积的迹在轮换乘法次序时不变的性质
\ech
\end{frame}

\begin{frame}
\chtitle{第1题解答}
\bch
$$\mathrm{Tr}(\slashed{A}\slashed{B}) = A_\mu B_\nu \mathrm{Tr}\left(\gamma^\mu\gamma^\nu\right)$$
因$\mathrm{Tr}(XY) = \mathrm{Tr}(YX)$，所以 
$$\mathrm{Tr}(\slashed{A}\slashed{B}) = A_\mu B_\nu \mathrm{Tr}\left(\gamma^\nu\gamma^\mu\right)$$
两式相加除以2，并利用$\gamma^\mu\gamma^\nu + \gamma^\nu\gamma^\mu = 2 g^{\mu\nu}I_{4\times 4}$，
$$\mathrm{Tr}(\slashed{A}\slashed{B}) = A_\mu B_\nu g^{\mu\nu} \mathrm{Tr}(I_{4\times 4}) = 4 A^\mu B_\mu$$
 
\ech
\end{frame}

\begin{frame}
\chtitle{第2题：题目和思路}
\bch
题目：对旋量$\psi$证明$\bar{\psi}\gamma^\mu\gamma^\nu\gamma^\rho\psi$为洛仑兹变换下的三阶张量。


\skipline
思路：这类题目总是要用$\Lambda^\dagger = \gamma^0\Lambda^{-1}\gamma^0$和$\Lambda$矩阵的定义
\ech
\end{frame}

\begin{frame}
\chtitle{第2题解答}
\bch
设洛仑兹变换$x^\mu \rightarrow a^\mu_{\ \nu} x^\nu$对应的旋量变换矩阵为$\Lambda$。利用课上证明了的$\Lambda^\dagger = \gamma^0\Lambda^{-1}\gamma^0$以及$\Lambda$矩阵的定义，
\bea
\bar{\psi}\gamma^\mu\gamma^\nu\gamma^\rho\psi &\rightarrow &\psi^\dagger\Lambda^\dagger \gamma^0 \gamma^\mu\gamma^\nu\gamma^\rho\Lambda\psi  \newl
&=&\psi^\dagger\gamma^0\Lambda^{-1}(\gamma^0)^2\gamma^\mu\gamma^\nu\gamma^\rho\Lambda\psi \newl
&=&\barpsi\Lambda^{-1}\gamma^\mu\gamma^\nu\gamma^\rho\Lambda\psi \newl
&=&\barpsi\Lambda^{-1}\gamma^\mu\Lambda\Lambda^{-1}\gamma^\nu\Lambda\Lambda^{-1}\gamma^\rho\Lambda\psi \newl
&=&\barpsi a^\mu_{\ \alpha}\gamma^\alpha a^\nu_{\ \beta}\gamma^\beta a^\rho_{\ \sigma}\gamma^\sigma \psi \newl
&=&(a^\mu_{\ \alpha} a^\nu_{\ \beta} a^\rho_{\ \sigma})\barpsi \gamma^\alpha \gamma^\beta \gamma^\sigma \psi
\eea
上式表明$\bar{\psi}\gamma^\mu\gamma^\nu\gamma^\rho\psi$满足三阶张量的定义。
\ech
\end{frame}

\begin{frame}
\chtitle{第3题：题目和思路}
\bch
题目：如果一个实标量场$\phi$和一个旋量场$\psi$有相互作用，拉氏密度为
$$\lagr = \frac{1}{2}\partial_\mu\phi\partial^\mu\phi - \frac{1}{2}m^2\phi^2 - \frac{g}{2}\phi^2\bar\psi\psi + \bar\psi(i\slashed{\partial}-m)\psi$$
其中$g$为耦合常数。

试推导$\phi$和$\psi$的运动方程。

\skipline
思路：对$\phi$和$\bar\psi$写出Euler-Lagrange方程即可
\ech
\end{frame}

\begin{frame}
\chtitle{第3题解答}
\bch
直接利用Euler-Lagrange方程得到$\phi$的运动方程
$$ (\partial^2 + m^2 + g\barpsi\psi) \phi = 0$$
$\psi$的运动方程(通过取$\barpsi$的Euler-Lagrange方程得到)
$$ (-\frac{g}{2}\phi^2 + i\slashed{\partial} - m) \psi = 0 $$
\ech
\end{frame}

\begin{frame}
\chtitle{第4题：题目和思路}
\bch
题目：设有三维动量$\veck = (k_x, k_y, k_z)$。请在以$z$方向自旋向上的态$|\uparrow\rangle$和自旋向下的态$|\downarrow\rangle$为基的表象里，写出沿$\veck$方向的电子自旋算符的矩阵表达式，并求它的所有本征值$s$和本征矢$\zeta_{\veck,s}$。

\skipline
思路：同一表象下的算符的和的矩阵表示等于算符的矩阵表示的和。

\ech
\end{frame}


\begin{frame}
\chtitle{第4题解答}
\bch
沿$\veck$方向的自旋算符为
\be
\frac{1}{2} \sigma\cdot\frac{\veck}{|\veck|} = \frac{k_x}{2|\veck|}\sigma^1 + \frac{k_y}{2|\veck|}\sigma^2 + \frac{k_z}{2|\veck|}\sigma^3 = 
\frac{1}{2|\veck|}
\bmat{cc}
k_z & k_x - \ii  k_y \\
k_x + \ii k_y & -k_z
\emat
\ee
本征值$s$满足：
$$(k_z-2|\veck|s)(-k_z - 2|\veck|s) - (k_x+\ii k_y)(k_x-\ii k_y) = 0$$
化简即得$s = \pm 1/2$
对应$s = \pm 1/2$的本征矢量分别为:
{\scriptsize
\be
\zeta_{\veck, 1/2} = \frac{1}{\sqrt{2|\veck|(|\veck|+k_z)}}
\bmat{c}
k_z+|\veck|\\
k_x+\ii k_y
\emat,\ 
\zeta_{\veck, -1/2} = \frac{1}{\sqrt{2|\veck|(|\veck|+k_z)}}
\bmat{c}
k_x-\ii k_y \\
-k_z-|\veck|
\emat
\ee
}

\ech
\end{frame}

\begin{frame}
\chtitle{第5题：题目和思路}
\bch
题目：对$m=0$的旋量和非零的三维动量$\veck$，记相应的四维动量为$k$，证明$\slashed{k}$只有两个线性独立的本征态。 

\skipline
思路：由于矩阵的本征态个数不依赖于表象，我们可以简单取$z$轴沿$\veck$方向写出$\slashed{k}$。
\ech
\end{frame}

\begin{frame}
\chtitle{第5题解答}
\bch
{\small
取$\veck$方向为$z$轴，则$k_\mu = (|\veck|, 0 , 0, -|\veck|)$}
{\scriptsize
\be
\slashed{k} = \gamma^0k_0 + \gamma^3k_3 =
|\veck|
\bmat{rrrr}
1 & 0 & -1 & 0 \\
0 & 1 & 0 & 1 \\
1 & 0 & -1 & 0 \\
0 & -1 & 0 & -1 
\emat
\ee
}
{\small
设$\slashed{k}$有本征值$\lambda$和非零本征矢$\psi$,则利用课上的结论$\slashed{k}^2=m^2=0$有
$$\lambda^2\psi = \lambda\slashed{k}\psi = \slashed{k}^2\psi  = 0$$ 
故$\lambda = 0$。
所以$\slashed{k}$的本征矢$\psi$满足}
{\scriptsize
\be
\bmat{rrrr}
1 & 0 & -1 & 0 \\
0 & 1 & 0 & 1 \\
1 & 0 & -1 & 0 \\
0 & -1 & 0 & -1 
\emat
\bmat{c}
\psi_1\\
\psi_2\\
\psi_3\\
\psi_4
\emat = 0
\ee
}{\small
显然它只有两个线性独立的解分别对应于$\psi_1 = \psi_3$和$\psi_2 = -\psi_4$。
}
\ech
\end{frame}

\begin{frame}
\chtitle{第5题解答}
\bch
{\small
或者更清晰的说法是, 上面的方程等价于$\psi_1 = \psi_3$, $\psi_2 = -\psi_4$，即解可以写成
\be
\psi = \psi_1 \bmat{c}
1 \\
0 \\
1 \\
0
\emat  + 
\psi_2 \bmat{c}
0 \\
1 \\
0 \\
-1
\emat
\ee
也就是说只有两个线性独立的本征矢。
}
\ech
\end{frame}

\end{document}
